\chapter{Introdução}
\pagestyle{plain}

A Inteligência Artificial (IA) consolidou-se como uma das tecnologias mais transformadoras da atualidade, promovendo mudanças profundas na forma como organizações desenvolvem produtos, prestam serviços, automatizam processos e utilizam dados para tomada de decisão. Nos últimos anos, a evolução dos modelos de Machine Learning, Deep Learning e Inteligência Artificial Generativa ampliou significativamente o potencial de aplicação dessas tecnologias em praticamente todos os setores da economia.
A crescente disponibilidade de grandes volumes de dados, aliada ao aumento da capacidade computacional e ao desenvolvimento de modelos de linguagem de larga escala (Large Language Models – LLMs), tornou possível automatizar atividades anteriormente restritas à atuação humana, como interpretação de documentos, desenvolvimento de software, geração de conteúdo, análise preditiva, inspeção visual, atendimento inteligente, automação de processos e apoio à decisão.

Paralelamente, a transformação digital da indústria impulsionou a integração entre Inteligência Artificial, Internet das Coisas (IoT), Computação em Nuvem, Computação de Borda (Edge Computing), Big Data e Computação Distribuída, formando o conjunto de tecnologias que sustenta a Indústria 4.0.

Nesse contexto, observa-se uma crescente demanda por profissionais capazes de desenvolver soluções inteligentes completas, compreendendo não apenas o treinamento de modelos de IA, mas também todo o ciclo de vida das soluções, incluindo coleta e tratamento de dados, engenharia de software, integração com sistemas corporativos, automação, implantação em produção, monitoramento contínuo, governança e evolução dos modelos.

Ao mesmo tempo, a rápida evolução da IA Generativa transformou significativamente a Engenharia de Software. Ferramentas como assistentes de programação, plataformas de automação inteligente, agentes autônomos e modelos multimodais passaram a fazer parte do cotidiano dos desenvolvedores, exigindo novas competências relacionadas à engenharia de prompts, engenharia de contexto, construção de agentes inteligentes, Retrieval-Augmented Generation (RAG), integração com APIs, observabilidade e MLOps.

Diante desse cenário, torna-se indispensável a formação de profissionais que compreendam de forma integrada os fundamentos científicos da Inteligência Artificial, as metodologias modernas de Ciência de Dados, as práticas contemporâneas de Engenharia de Software e as tecnologias utilizadas na implantação de soluções inteligentes em ambientes corporativos e industriais.

O projeto foi desenvolvido no intuito de atender a essa necessidade, estruturando uma formação abrangente que integrará conhecimentos de Ciência de Dados, Machine Learning, Deep Learning, Inteligência Artificial Generativa, Desenvolvimento de Software Assistido por IA, Engenharia de Prompt, Agentes Inteligentes, Visão Computacional, Internet das Coisas (IoT), Computação Distribuída, MLOps e DevOps para Inteligência Artificial.

Sua organização curricular foi desenvolvida segundo uma progressão pedagógica baseada em competências, o que permitirá que o estudante evolua desde os fundamentos matemáticos e computacionais até o desenvolvimento de arquiteturas avançadas de Inteligência Artificial aplicadas a problemas reais de software, indústria e negócios.

A proposta pedagógica privilegia a aprendizagem baseada em projetos, resolução de problemas, experimentação prática e desenvolvimento de soluções completas assim como prevê a Matodologia Senai de Educação Profissional (MSEP), aproximando o estudante das tecnologias, metodologias e ferramentas que serão (ou são atualmente) utilizadas por empresas de tecnologia, startups, centros de pesquisa e organizações industriais.